\documentclass[11pt,a4paper]{article}
\usepackage[utf8]{inputenc}
\usepackage[T1]{fontenc}
\usepackage[french]{babel}
\usepackage{amsmath}
\usepackage{amssymb}
\usepackage{amsthm}
\usepackage{mathtools}
\usepackage{geometry}
\usepackage{enumitem}
\usepackage{xcolor}
\usepackage{titlesec}
\usepackage{fancyhdr}
\usepackage{booktabs}
\usepackage{array}
\usepackage{longtable}
\usepackage{tcolorbox}
\usepackage{multicol}

% Configuration de la page
\geometry{margin=2.5cm}
\pagestyle{fancy}
\fancyhf{}
\fancyhead[L]{\leftmark}
\fancyhead[R]{Exercices Lois de Probabilité}
\fancyfoot[C]{\thepage}

% Configuration des titres
\titleformat{\section}
{\Large\bfseries\color{blue!70!black}}
{}
{0em}
{}[\titlerule]

\titleformat{\subsection}
{\large\bfseries\color{blue!50!black}}
{}
{0em}
{}

\titleformat{\subsubsection}
{\normalsize\bfseries\color{blue!30!black}}
{}
{0em}
{}

% Boîtes colorées
\newtcolorbox{astucebox}{
    colback=green!5!white,
    colframe=green!75!black,
    title=\textbf{Astuce},
    fonttitle=\bfseries
}

\newtcolorbox{rappelbox}{
    colback=blue!5!white,
    colframe=blue!75!black,
    title=\textbf{Rappel},
    fonttitle=\bfseries
}

\newtcolorbox{attentionbox}{
    colback=orange!5!white,
    colframe=orange!75!black,
    title=\textbf{Attention},
    fonttitle=\bfseries
}

\newtcolorbox{exemplebox}{
    colback=yellow!5!white,
    colframe=yellow!75!black,
    title=\textbf{Exemple},
    fonttitle=\bfseries
}

% Commandes personnalisées
\newcommand{\R}{\mathbb{R}}
\newcommand{\Var}{\text{Var}}
\newcommand{\E}{\mathbb{E}}
\newcommand{\N}{\mathcal{N}}
\newcommand{\Bin}{\text{Bin}}
\newcommand{\Exp}{\text{Exp}}
\newcommand{\Ber}{\text{Ber}}

% Métadonnées
\title{Exercices Complets - Lois de Probabilité\\
\large Binomiale, Bernoulli, Exponentielle, Normale, Normale Centrée Réduite}
\author{AmineGR03}
\date{\today}

\begin{document}

\maketitle

\tableofcontents
\newpage

\section{Introduction}

Ce document présente des exercices pratiques sur les principales lois de probabilité :
\begin{itemize}
    \item Loi de Bernoulli
    \item Loi binomiale
    \item Loi exponentielle
    \item Loi normale
    \item Loi normale centrée réduite
\end{itemize}

Chaque exercice comprend :
\begin{itemize}
    \item Un énoncé détaillé
    \item Des rappels théoriques
    \item Une solution complète avec explications étape par étape
    \item Des astuces et pièges à éviter
\end{itemize}

\newpage

\section{Exercice 1 : Loi de Bernoulli}

\subsection{Énoncé}

Une machine produit des pièces. La probabilité qu'une pièce soit défectueuse est $p = 0.15$. On prélève une pièce au hasard.

\begin{enumerate}
    \item Quelle est la loi de probabilité du nombre de pièces défectueuses ?
    \item Calculer l'espérance et la variance.
    \item Quelle est la probabilité que la pièce soit défectueuse ?
    \item Quelle est la probabilité que la pièce soit conforme ?
\end{enumerate}

\subsection{Rappels Théoriques}

\begin{rappelbox}
\textbf{Loi de Bernoulli $\Ber(p)$ :}

Une variable aléatoire $X$ suit une loi de Bernoulli de paramètre $p$ si :
\begin{itemize}
    \item $X$ prend les valeurs 0 (échec) ou 1 (succès)
    \item $P(X = 1) = p$ (probabilité de succès)
    \item $P(X = 0) = 1 - p = q$ (probabilité d'échec)
\end{itemize}

\textbf{Propriétés :}
\begin{align}
\E(X) &= p \\
\Var(X) &= p(1-p) = pq
\end{align}
\end{rappelbox}

\subsection{Solution Détaillée}

\begin{enumerate}
    \item \textbf{Loi de probabilité :}
    
    Soit $X$ la variable aléatoire qui vaut 1 si la pièce est défectueuse, et 0 sinon.
    
    $X$ suit une loi de Bernoulli de paramètre $p = 0.15$ : $X \sim \Ber(0.15)$.
    
    \begin{align}
    P(X = 0) &= 1 - 0.15 = 0.85 \quad \text{(pièce conforme)} \\
    P(X = 1) &= 0.15 \quad \text{(pièce défectueuse)}
    \end{align}
    
    \item \textbf{Espérance et variance :}
    
    \begin{align}
    \E(X) &= p = 0.15 \\
    \Var(X) &= p(1-p) = 0.15 \times 0.85 = 0.1275
    \end{align}
    
    \item \textbf{Probabilité que la pièce soit défectueuse :}
    
    $$P(X = 1) = p = 0.15 = 15\%$$
    
    \item \textbf{Probabilité que la pièce soit conforme :}
    
    $$P(X = 0) = 1 - p = 0.85 = 85\%$$
\end{enumerate}

\begin{astucebox}
La loi de Bernoulli est la base de nombreuses autres lois (binomiale, géométrique). Elle modélise une expérience avec deux issues possibles : succès ou échec.
\end{astucebox}

\newpage

\section{Exercice 2 : Loi Binomiale}

\subsection{Énoncé}

On prélève 20 pièces au hasard dans la production. La probabilité qu'une pièce soit défectueuse est $p = 0.15$.

\begin{enumerate}
    \item Quelle est la loi de probabilité du nombre de pièces défectueuses ?
    \item Calculer l'espérance et la variance.
    \item Quelle est la probabilité d'avoir exactement 3 pièces défectueuses ?
    \item Quelle est la probabilité d'avoir au moins 2 pièces défectueuses ?
    \item Quelle est la probabilité d'avoir au plus 1 pièce défectueuse ?
\end{enumerate}

\subsection{Rappels Théoriques}

\begin{rappelbox}
\textbf{Loi binomiale $\Bin(n, p)$ :}

Une variable aléatoire $X$ suit une loi binomiale de paramètres $n$ et $p$ si :
\begin{itemize}
    \item $X$ compte le nombre de succès dans $n$ épreuves de Bernoulli indépendantes
    \item Chaque épreuve a une probabilité de succès $p$
\end{itemize}

\textbf{Fonction de masse :}
$$P(X = k) = \binom{n}{k} p^k (1-p)^{n-k} = \binom{n}{k} p^k q^{n-k}$$

où $\binom{n}{k} = \frac{n!}{k!(n-k)!}$ est le coefficient binomial.

\textbf{Propriétés :}
\begin{align}
\E(X) &= np \\
\Var(X) &= np(1-p) = npq
\end{align}
\end{rappelbox}

\subsection{Solution Détaillée}

\begin{enumerate}
    \item \textbf{Loi de probabilité :}
    
    Soit $X$ le nombre de pièces défectueuses parmi les 20.
    
    $X$ suit une loi binomiale de paramètres $n = 20$ et $p = 0.15$ : $X \sim \Bin(20, 0.15)$.
    
    \item \textbf{Espérance et variance :}
    
    \begin{align}
    \E(X) &= np = 20 \times 0.15 = 3 \\
    \Var(X) &= np(1-p) = 20 \times 0.15 \times 0.85 = 2.55
    \end{align}
    
    \item \textbf{Probabilité d'avoir exactement 3 pièces défectueuses :}
    
    $$P(X = 3) = \binom{20}{3} (0.15)^3 (0.85)^{17}$$
    
    Calculons :
    \begin{align}
    \binom{20}{3} &= \frac{20!}{3! \times 17!} = \frac{20 \times 19 \times 18}{3 \times 2 \times 1} = \frac{6840}{6} = 1140 \\
    P(X = 3) &= 1140 \times (0.15)^3 \times (0.85)^{17} \\
    &= 1140 \times 0.003375 \times 0.0631 \\
    &\approx 1140 \times 0.000213 \\
    &\approx 0.243
    \end{align}
    
    Donc $P(X = 3) \approx 0.243 = 24.3\%$.
    
    \item \textbf{Probabilité d'avoir au moins 2 pièces défectueuses :}
    
    $$P(X \geq 2) = 1 - P(X < 2) = 1 - [P(X = 0) + P(X = 1)]$$
    
    Calculons d'abord :
    \begin{align}
    P(X = 0) &= \binom{20}{0} (0.15)^0 (0.85)^{20} = 1 \times 1 \times (0.85)^{20} \approx 0.0388 \\
    P(X = 1) &= \binom{20}{1} (0.15)^1 (0.85)^{19} = 20 \times 0.15 \times (0.85)^{19} \approx 0.137
    \end{align}
    
    Donc :
    $$P(X \geq 2) = 1 - (0.0388 + 0.137) = 1 - 0.1758 = 0.8242 \approx 82.4\%$$
    
    \item \textbf{Probabilité d'avoir au plus 1 pièce défectueuse :}
    
    $$P(X \leq 1) = P(X = 0) + P(X = 1) = 0.0388 + 0.137 = 0.1758 \approx 17.6\%$$
\end{enumerate}

\begin{attentionbox}
Pour calculer $\binom{n}{k}$, on peut utiliser la formule :
$$\binom{n}{k} = \frac{n \times (n-1) \times \ldots \times (n-k+1)}{k \times (k-1) \times \ldots \times 1}$$
\end{attentionbox}

\newpage

\section{Exercice 3 : Loi Exponentielle}

\subsection{Énoncé}

Le temps d'attente (en minutes) avant l'arrivée d'un bus suit une loi exponentielle de paramètre $\lambda = 0.2$ (bus/minute).

\begin{enumerate}
    \item Quelle est la probabilité d'attendre moins de 5 minutes ?
    \item Quelle est la probabilité d'attendre entre 3 et 7 minutes ?
    \item Quelle est la probabilité d'attendre plus de 10 minutes ?
    \item Calculer l'espérance et la variance.
    \item Quel est le temps médian d'attente ?
\end{enumerate}

\subsection{Rappels Théoriques}

\begin{rappelbox}
\textbf{Loi exponentielle $\Exp(\lambda)$ :}

Une variable aléatoire $X$ suit une loi exponentielle de paramètre $\lambda > 0$ si sa densité est :
$$f(x) = \begin{cases}
\lambda e^{-\lambda x} & \text{si } x \geq 0 \\
0 & \text{si } x < 0
\end{cases}$$

\textbf{Fonction de répartition :}
$$F(x) = P(X \leq x) = \begin{cases}
1 - e^{-\lambda x} & \text{si } x \geq 0 \\
0 & \text{si } x < 0
\end{cases}$$

\textbf{Propriétés :}
\begin{align}
\E(X) &= \frac{1}{\lambda} \\
\Var(X) &= \frac{1}{\lambda^2}
\end{align}

\textbf{Propriété d'absence de mémoire :}
$$P(X > s + t | X > s) = P(X > t)$$
\end{rappelbox}

\subsection{Solution Détaillée}

\begin{enumerate}
    \item \textbf{Probabilité d'attendre moins de 5 minutes :}
    
    $$P(X < 5) = F(5) = 1 - e^{-\lambda \times 5} = 1 - e^{-0.2 \times 5} = 1 - e^{-1}$$
    
    $$P(X < 5) = 1 - e^{-1} \approx 1 - 0.3679 = 0.6321 \approx 63.2\%$$
    
    \item \textbf{Probabilité d'attendre entre 3 et 7 minutes :}
    
    $$P(3 < X < 7) = F(7) - F(3) = (1 - e^{-0.2 \times 7}) - (1 - e^{-0.2 \times 3})$$
    
    $$P(3 < X < 7) = e^{-0.6} - e^{-1.4} = 0.5488 - 0.2466 = 0.3022 \approx 30.2\%$$
    
    \item \textbf{Probabilité d'attendre plus de 10 minutes :}
    
    $$P(X > 10) = 1 - F(10) = 1 - (1 - e^{-0.2 \times 10}) = e^{-2} \approx 0.1353 \approx 13.5\%$$
    
    \item \textbf{Espérance et variance :}
    
    \begin{align}
    \E(X) &= \frac{1}{\lambda} = \frac{1}{0.2} = 5 \text{ minutes} \\
    \Var(X) &= \frac{1}{\lambda^2} = \frac{1}{0.04} = 25 \text{ minutes}^2
    \end{align}
    
    \item \textbf{Temps médian :}
    
    Le temps médian $m$ vérifie $P(X \leq m) = 0.5$, donc :
    $$F(m) = 1 - e^{-\lambda m} = 0.5$$
    
    $$e^{-\lambda m} = 0.5$$
    $$-\lambda m = \ln(0.5) = -\ln(2)$$
    $$m = \frac{\ln(2)}{\lambda} = \frac{\ln(2)}{0.2} = \frac{0.693}{0.2} = 3.465 \text{ minutes}$$
    
    Le temps médian est d'environ 3.47 minutes.
\end{enumerate}

\begin{astucebox}
La loi exponentielle modélise souvent les temps d'attente. Elle possède la propriété d'absence de mémoire : peu importe le temps déjà attendu, la probabilité d'attendre encore $t$ minutes est la même.
\end{astucebox}

\newpage

\section{Exercice 4 : Loi Normale}

\subsection{Énoncé}

Le poids des étudiants d'une université suit une loi normale de moyenne $\mu = 70$ kg et d'écart-type $\sigma = 8$ kg.

\begin{enumerate}
    \item Quelle est la probabilité qu'un étudiant pèse moins de 65 kg ?
    \item Quelle est la probabilité qu'un étudiant pèse entre 65 et 75 kg ?
    \item Quelle est la probabilité qu'un étudiant pèse plus de 80 kg ?
    \item Quel est le poids tel que 90\% des étudiants pèsent moins ?
    \item Calculer l'espérance et la variance.
\end{enumerate}

\subsection{Rappels Théoriques}

\begin{rappelbox}
\textbf{Loi normale $\mathcal{N}(\mu, \sigma^2)$ :}

Une variable aléatoire $X$ suit une loi normale de paramètres $\mu$ et $\sigma^2$ si sa densité est :
$$f(x) = \frac{1}{\sigma\sqrt{2\pi}} e^{-\frac{1}{2}\left(\frac{x-\mu}{\sigma}\right)^2}$$

\textbf{Propriétés :}
\begin{itemize}
    \item Symétrique autour de $\mu$
    \item $\E(X) = \mu$
    \item $\Var(X) = \sigma^2$
    \item $\sigma(X) = \sigma$ (écart-type)
\end{itemize}

\textbf{Standardisation :}
Si $X \sim \mathcal{N}(\mu, \sigma^2)$, alors :
$$Z = \frac{X - \mu}{\sigma} \sim \mathcal{N}(0, 1)$$
\end{rappelbox}

\subsection{Solution Détaillée}

Soit $X \sim \mathcal{N}(70, 8^2)$ le poids d'un étudiant.

\begin{enumerate}
    \item \textbf{Probabilité qu'un étudiant pèse moins de 65 kg :}
    
    On standardise : $Z = \frac{X - 70}{8}$
    
    $$P(X < 65) = P\left(Z < \frac{65 - 70}{8}\right) = P(Z < -0.625)$$
    
    Par symétrie : $P(Z < -0.625) = P(Z > 0.625) = 1 - \Phi(0.625)$
    
    En consultant la table de la loi normale centrée réduite :
    $$\Phi(0.625) \approx 0.734$$
    
    Donc : $P(X < 65) = 1 - 0.734 = 0.266 \approx 26.6\%$
    
    \item \textbf{Probabilité qu'un étudiant pèse entre 65 et 75 kg :}
    
    $$P(65 < X < 75) = P\left(\frac{65 - 70}{8} < Z < \frac{75 - 70}{8}\right) = P(-0.625 < Z < 0.625)$$
    
    $$P(65 < X < 75) = \Phi(0.625) - \Phi(-0.625) = 0.734 - 0.266 = 0.468 \approx 46.8\%$$
    
    \item \textbf{Probabilité qu'un étudiant pèse plus de 80 kg :}
    
    $$P(X > 80) = P\left(Z > \frac{80 - 70}{8}\right) = P(Z > 1.25)$$
    
    $$P(X > 80) = 1 - \Phi(1.25) = 1 - 0.8944 = 0.1056 \approx 10.6\%$$
    
    \item \textbf{Poids tel que 90\% des étudiants pèsent moins :}
    
    On cherche $x$ tel que $P(X < x) = 0.90$.
    
    $$P(X < x) = P\left(Z < \frac{x - 70}{8}\right) = 0.90$$
    
    En consultant la table : $P(Z < 1.28) \approx 0.90$
    
    Donc : $\frac{x - 70}{8} = 1.28$, soit $x = 70 + 1.28 \times 8 = 70 + 10.24 = 80.24$ kg
    
    90\% des étudiants pèsent moins de 80.24 kg.
    
    \item \textbf{Espérance et variance :}
    
    \begin{align}
    \E(X) &= \mu = 70 \text{ kg} \\
    \Var(X) &= \sigma^2 = 8^2 = 64 \text{ kg}^2
    \end{align}
\end{enumerate}

\newpage

\section{Exercice 5 : Loi Normale Centrée Réduite}

\subsection{Énoncé}

Soit $Z \sim \mathcal{N}(0, 1)$ une variable aléatoire suivant une loi normale centrée réduite.

\begin{enumerate}
    \item Quelle est la probabilité que $Z$ soit inférieur à 1.5 ?
    \item Quelle est la probabilité que $Z$ soit entre -1 et 1 ?
    \item Quelle est la probabilité que $Z$ soit supérieur à 2 ?
    \item Quelle est la valeur $z$ telle que $P(Z < z) = 0.95$ ?
    \item Quelle est la valeur $z$ telle que $P(|Z| < z) = 0.90$ ?
    \item Calculer l'espérance et la variance.
\end{enumerate}

\subsection{Rappels Théoriques}

\begin{rappelbox}
\textbf{Loi normale centrée réduite $\mathcal{N}(0, 1)$ :}

C'est une loi normale avec $\mu = 0$ et $\sigma = 1$.

\textbf{Densité :}
$$\varphi(z) = \frac{1}{\sqrt{2\pi}} e^{-\frac{z^2}{2}}$$

\textbf{Fonction de répartition :}
$$\Phi(z) = P(Z \leq z) = \int_{-\infty}^{z} \varphi(t) dt$$

\textbf{Propriétés :}
\begin{itemize}
    \item Symétrique : $\Phi(-z) = 1 - \Phi(z)$
    \item $\E(Z) = 0$
    \item $\Var(Z) = 1$
    \item $\sigma(Z) = 1$
\end{itemize}

\textbf{Valeurs remarquables :}
\begin{itemize}
    \item $P(|Z| < 1) \approx 0.68$ (68\%)
    \item $P(|Z| < 2) \approx 0.95$ (95\%)
    \item $P(|Z| < 3) \approx 0.997$ (99.7\%)
\end{itemize}
\end{rappelbox}

\subsection{Solution Détaillée}

\begin{enumerate}
    \item \textbf{Probabilité que $Z$ soit inférieur à 1.5 :}
    
    $$P(Z < 1.5) = \Phi(1.5)$$
    
    En consultant la table : $\Phi(1.5) \approx 0.9332$
    
    Donc $P(Z < 1.5) \approx 0.9332 = 93.32\%$
    
    \item \textbf{Probabilité que $Z$ soit entre -1 et 1 :}
    
    $$P(-1 < Z < 1) = \Phi(1) - \Phi(-1) = \Phi(1) - (1 - \Phi(1)) = 2\Phi(1) - 1$$
    
    En consultant la table : $\Phi(1) \approx 0.8413$
    
    $$P(-1 < Z < 1) = 2 \times 0.8413 - 1 = 1.6826 - 1 = 0.6826 \approx 68.3\%$$
    
    \item \textbf{Probabilité que $Z$ soit supérieur à 2 :}
    
    $$P(Z > 2) = 1 - \Phi(2)$$
    
    En consultant la table : $\Phi(2) \approx 0.9772$
    
    $$P(Z > 2) = 1 - 0.9772 = 0.0228 \approx 2.28\%$$
    
    \item \textbf{Valeur $z$ telle que $P(Z < z) = 0.95$ :}
    
    On cherche le quantile d'ordre 0.95.
    
    En consultant la table inverse : $z_{0.95} \approx 1.645$
    
    Donc $P(Z < 1.645) \approx 0.95$
    
    \item \textbf{Valeur $z$ telle que $P(|Z| < z) = 0.90$ :}
    
    $$P(|Z| < z) = P(-z < Z < z) = 2\Phi(z) - 1 = 0.90$$
    
    Donc : $2\Phi(z) = 1.90$, soit $\Phi(z) = 0.95$
    
    En consultant la table inverse : $z \approx 1.645$
    
    Donc $P(|Z| < 1.645) \approx 0.90$
    
    \item \textbf{Espérance et variance :}
    
    \begin{align}
    \E(Z) &= 0 \\
    \Var(Z) &= 1
    \end{align}
\end{enumerate}

\begin{astucebox}
La loi normale centrée réduite est fondamentale car toute loi normale peut être standardisée. Les tables de la loi normale centrée réduite permettent de calculer les probabilités pour n'importe quelle loi normale.
\end{astucebox}

\newpage

\section{Résumé et Tableaux Récapitulatifs}

\subsection{Tableau Comparatif des Lois}

\begin{table}[h]
\centering
\small
\begin{tabular}{|l|c|c|c|}
\hline
\textbf{Loi} & \textbf{Espérance} & \textbf{Variance} & \textbf{Paramètres} \\
\hline
Bernoulli $\Ber(p)$ & $p$ & $p(1-p)$ & $p \in [0,1]$ \\
\hline
Binomiale $\Bin(n,p)$ & $np$ & $np(1-p)$ & $n \in \mathbb{N}^*, p \in [0,1]$ \\
\hline
Exponentielle $\Exp(\lambda)$ & $\frac{1}{\lambda}$ & $\frac{1}{\lambda^2}$ & $\lambda > 0$ \\
\hline
Normale $\mathcal{N}(\mu, \sigma^2)$ & $\mu$ & $\sigma^2$ & $\mu \in \mathbb{R}, \sigma > 0$ \\
\hline
Normale centrée réduite $\mathcal{N}(0,1)$ & $0$ & $1$ & Aucun \\
\hline
\end{tabular}
\caption{Propriétés des principales lois de probabilité}
\end{table}

\subsection{Formules Clés}

\begin{rappelbox}
\textbf{Loi de Bernoulli :}
$$P(X = k) = \begin{cases}
p & \text{si } k = 1 \\
1-p & \text{si } k = 0
\end{cases}$$

\textbf{Loi binomiale :}
$$P(X = k) = \binom{n}{k} p^k (1-p)^{n-k}, \quad k = 0, 1, \ldots, n$$

\textbf{Loi exponentielle :}
$$F(x) = 1 - e^{-\lambda x}, \quad x \geq 0$$

\textbf{Loi normale :}
$$P(a < X < b) = \Phi\left(\frac{b-\mu}{\sigma}\right) - \Phi\left(\frac{a-\mu}{\sigma}\right)$$

\textbf{Standardisation :}
$$Z = \frac{X - \mu}{\sigma} \sim \mathcal{N}(0, 1)$$
\end{rappelbox}

\vspace{2cm}

\begin{center}
\textbf{Fin du document}
\end{center}

\end{document}



